\section{Khái Niệm Mất Cân Bằng Dữ Liệu}
% ============================================================

Mất cân bằng dữ liệu (\textit{data imbalance} hay \textit{class imbalance}) là hiện tượng phân phối các lớp trong tập dữ liệu phân loại không đồng đều một cách đáng kể, trong đó một hoặc nhiều lớp chiếm tỉ lệ áp đảo so với các lớp còn lại. Lớp có số lượng mẫu lớn hơn được gọi là lớp đa số (\textit{majority class}), trong khi lớp có số lượng mẫu ít hơn được gọi là lớp thiểu số (\textit{minority class}). Tỉ lệ mất cân bằng (\textit{imbalance ratio} --- IR) thường được đo bằng:

\begin{equation}
  \text{IR} = \frac{|\text{majority class}|}{|\text{minority class}|}
  \label{eq:ir}
\end{equation}

Trong thực tế, tỉ lệ này có thể dao động từ 10:1 (mất cân bằng nhẹ) đến 1000:1 hoặc cao hơn (mất cân bằng nghiêm trọng). Chẳng hạn, trong bài toán phát hiện gian lận thẻ tín dụng, chỉ khoảng $0{,}1\%$--$0{,}5\%$ giao dịch là gian lận, tức là tỉ lệ mất cân bằng có thể lên đến 200:1 hay hơn.

% ============================================================
\section{Nguyên Nhân Dẫn Đến Mất Cân Bằng Dữ Liệu}
% ============================================================

Hiện tượng mất cân bằng dữ liệu phát sinh từ nhiều nguyên nhân khác nhau, phản ánh bản chất tự nhiên của các sự kiện trong thế giới thực. Thứ nhất, sự hiếm gặp tự nhiên (\textit{natural rarity}) là nguyên nhân phổ biến nhất: các sự kiện bất thường như gian lận tài chính, chẩn đoán bệnh hiếm gặp hay sự cố kỹ thuật vốn dĩ xảy ra với tần suất rất thấp trong tự nhiên. Thứ hai, chi phí thu thập dữ liệu cao (\textit{high collection cost}) trong một số lĩnh vực --- chẳng hạn dữ liệu y tế từ bệnh nhân hiếm --- khiến việc thu thập đủ mẫu cho lớp thiểu số trở nên tốn kém hoặc bất khả thi. Thứ ba, sai lệch trong quá trình thu thập và gán nhãn dữ liệu (\textit{sampling bias}) cũng có thể tạo ra hoặc khuếch đại sự mất cân bằng một cách nhân tạo. Hiểu rõ nguyên nhân là bước đầu tiên quan trọng để lựa chọn chiến lược xử lý phù hợp.

% ============================================================
\section{Tác Động Đến Hiệu Suất Mô Hình}
% ============================================================

Dữ liệu mất cân bằng tạo ra những thách thức nghiêm trọng cho các thuật toán học máy truyền thống, vốn được thiết kế dựa trên giả định phân phối lớp đồng đều. Khi huấn luyện trên dữ liệu mất cân bằng, mô hình có xu hướng bị thiên lệch mạnh về phía lớp đa số (\textit{majority class bias}), dẫn đến việc phân loại hầu hết các mẫu vào lớp đa số để tối thiểu hóa hàm mất mát tổng thể. Hệ quả là mô hình đạt được độ chính xác tổng thể (\textit{accuracy}) cao nhưng hoàn toàn thất bại trong việc nhận diện lớp thiểu số --- chính là lớp thường có giá trị thực tiễn cao nhất trong bài toán. Ngoài ra, mất cân bằng dữ liệu còn làm giảm khả năng học của mô hình đối với các đặc trưng phân biệt của lớp thiểu số, dẫn đến ranh giới quyết định bị dịch chuyển về phía không thuận lợi cho lớp này.

% ============================================================
\section{Tại Sao Accuracy Không Còn Phù Hợp}
% ============================================================

Xét một bài toán nhị phân với tỉ lệ 99:1 giữa lớp đa số và lớp thiểu số. Một mô hình ngây thơ luôn dự đoán lớp đa số sẽ đạt accuracy $99\%$ mà không học được bất kỳ điều gì. Điều này minh họa rõ ràng rằng accuracy là một chỉ số đánh giá không đáng tin cậy trong bối cảnh mất cân bằng dữ liệu. Về mặt toán học, accuracy được định nghĩa là:

\begin{equation}
  \text{Accuracy} = \frac{TP + TN}{TP + TN + FP + FN}
  \label{eq:accuracy}
\end{equation}

\noindent trong đó $TP$, $TN$, $FP$, $FN$ lần lượt là số lượng dương tính thật (\textit{true positive}), âm tính thật (\textit{true negative}), dương tính giả (\textit{false positive}) và âm tính giả (\textit{false negative}). Khi lớp đa số áp đảo, $TN$ rất lớn sẽ đẩy accuracy lên cao bất kể $TP$ có giá trị như thế nào. Do đó, việc đánh giá mô hình trong bối cảnh mất cân bằng đòi hỏi các chỉ số tinh tế hơn, phản ánh trực tiếp hiệu năng trên lớp thiểu số.

% ============================================================
\section{Các Chỉ Số Đánh Giá Phù Hợp}
% ============================================================

\textbf{Precision} đo lường tỉ lệ các dự đoán dương tính thực sự là đúng, phản ánh chi phí của dương tính giả:

\begin{equation}
  \text{Precision} = \frac{TP}{TP + FP}
  \label{eq:precision}
\end{equation}

\textbf{Recall} (hay \textit{Sensitivity}) đo lường tỉ lệ các mẫu dương tính thực tế được phát hiện đúng, phản ánh chi phí của âm tính giả:

\begin{equation}
  \text{Recall} = \frac{TP}{TP + FN}
  \label{eq:recall}
\end{equation}

\textbf{F1-Score} là trung bình điều hòa của Precision và Recall, cung cấp một chỉ số cân bằng duy nhất khi cả hai đều quan trọng:

\begin{equation}
  \text{F1} = 2 \cdot \frac{\text{Precision} \times \text{Recall}}{\text{Precision} + \text{Recall}}
  \label{eq:f1}
\end{equation}

\textbf{ROC-AUC} (\textit{Area Under the Receiver Operating Characteristic Curve}) đo lường khả năng phân biệt tổng thể của mô hình trên tất cả các ngưỡng phân loại có thể. Giá trị AUC $= 1$ biểu thị phân loại hoàn hảo, trong khi AUC $= 0{,}5$ tương đương với phân loại ngẫu nhiên. Tuy nhiên, trong trường hợp mất cân bằng nghiêm trọng, ROC-AUC có thể vẫn cho giá trị cao dù hiệu năng trên lớp thiểu số kém.

\textbf{PR-AUC} (\textit{Area Under the Precision-Recall Curve}) được xem là chỉ số phù hợp hơn trong bối cảnh mất cân bằng nghiêm trọng, vì nó tập trung hoàn toàn vào hiệu năng trên lớp dương tính (lớp thiểu số) mà không bị ảnh hưởng bởi số lượng âm tính thật lớn.

% ============================================================
\section{Phân Loại Các Phương Pháp Xử Lý}
% ============================================================

Các phương pháp xử lý mất cân bằng dữ liệu được phân thành ba nhóm chính. Nhóm thứ nhất là \textbf{Oversampling} --- tăng số lượng mẫu của lớp thiểu số bằng cách tạo thêm dữ liệu. Nhóm thứ hai là \textbf{Undersampling} --- giảm số lượng mẫu của lớp đa số bằng cách loại bỏ một phần dữ liệu. Nhóm thứ ba là \textbf{Hybrid Methods} --- kết hợp cả hai chiến lược trên để tận dụng ưu điểm của từng phương pháp. Sự lựa chọn giữa các nhóm phụ thuộc vào kích thước tập dữ liệu, mức độ mất cân bằng và yêu cầu cụ thể của bài toán.

% ============================================================
\section{Trình Bày Chi Tiết Các Kỹ Thuật}
% ============================================================

% ------------------------------------------------------------
\subsection{Oversampling Methods}
% ------------------------------------------------------------

\subsubsection{Random Oversampling (ROS)}

\textbf{Khái niệm và nguyên lý.}
Random Oversampling là kỹ thuật đơn giản nhất trong nhóm oversampling, trong đó các mẫu từ lớp thiểu số được lựa chọn ngẫu nhiên và nhân bản (với thay thế) cho đến khi đạt được tỉ lệ lớp mong muốn. Về quy trình, thuật toán lặp lại việc lấy mẫu ngẫu nhiên từ tập lớp thiểu số và thêm các bản sao vào tập huấn luyện:

\begin{equation}
  \mathcal{D}'_{\min}
  = \mathcal{D}_{\min}
    \cup \bigl\{\mathbf{x} \mid \mathbf{x}
    \sim \mathrm{Uniform}(\mathcal{D}_{\min})\bigr\}
  \label{eq:ros}
\end{equation}

\begin{algorithm}[H]
\caption{Random Oversampling (ROS)}
\begin{algorithmic}[1]
  \Require Tập dữ liệu $\mathcal{D}$, tỉ lệ mục tiêu $r$
  \Ensure Tập dữ liệu cân bằng $\mathcal{D}'$
  \State $\mathcal{D}_{\min} \leftarrow$ các mẫu thuộc lớp thiểu số trong $\mathcal{D}$
  \State $N_{\text{need}} \leftarrow r \cdot |\mathcal{D}_{\text{maj}}| - |\mathcal{D}_{\min}|$
  \For{$i = 1$ \textbf{to} $N_{\text{need}}$}
    \State Chọn ngẫu nhiên $\mathbf{x} \sim \mathrm{Uniform}(\mathcal{D}_{\min})$
    \State $\mathcal{D}_{\min} \leftarrow \mathcal{D}_{\min} \cup \{\mathbf{x}\}$
  \EndFor
  \State \Return $\mathcal{D}' = \mathcal{D}_{\text{maj}} \cup \mathcal{D}_{\min}$
\end{algorithmic}
\end{algorithm}

\textbf{Ưu điểm:} Đơn giản, dễ triển khai, không làm mất thông tin, và không tạo ra dữ liệu tổng hợp nên không có rủi ro phân phối giả tạo.

\textbf{Nhược điểm:} Do chỉ nhân bản các mẫu hiện có mà không tạo ra thông tin mới, ROS dễ dẫn đến overfitting nghiêm trọng --- mô hình có thể ``ghi nhớ'' các mẫu thiểu số được lặp lại thay vì học các đặc trưng tổng quát hóa.

\textbf{Trường hợp nên sử dụng:} Phù hợp khi tập dữ liệu nhỏ, tỉ lệ mất cân bằng không quá nghiêm trọng, và khi cần một baseline nhanh trước khi thử các kỹ thuật phức tạp hơn.

% ------------------------------------------------------------
\subsubsection{SMOTE (Synthetic Minority Oversampling Technique)}

\textbf{Khái niệm.}
SMOTE, được đề xuất bởi Chawla et al. (2002), là kỹ thuật oversampling nền tảng và có ảnh hưởng nhất trong lĩnh vực xử lý mất cân bằng dữ liệu. Thay vì nhân bản các mẫu hiện có, SMOTE tạo ra các mẫu tổng hợp mới bằng cách nội suy trong không gian đặc trưng giữa các mẫu thiểu số lân cận.

\textbf{Nguyên lý hoạt động.}
Với mỗi mẫu $\mathbf{x}_i$ trong lớp thiểu số, SMOTE xác định $k$ láng giềng gần nhất (\textit{k-nearest neighbors} --- kNN) trong cùng lớp. Một mẫu tổng hợp mới $\mathbf{x}_{\text{new}}$ được tạo ra bằng công thức nội suy tuyến tính:

\begin{equation}
  \mathbf{x}_{\text{new}}
  = \mathbf{x}_i + \lambda \cdot (\mathbf{x}_{nn} - \mathbf{x}_i),
  \quad \lambda \sim \mathrm{Uniform}(0,\,1)
  \label{eq:smote}
\end{equation}

\noindent trong đó $\mathbf{x}_{nn}$ là một láng giềng được chọn ngẫu nhiên từ $k$ láng giềng gần nhất của $\mathbf{x}_i$.

\begin{algorithm}[H]
\caption{SMOTE}
\begin{algorithmic}[1]
  \Require $\mathcal{D}_{\min}$, số láng giềng $k$, số mẫu cần tạo $N$
  \Ensure Tập mẫu tổng hợp $\mathcal{S}$
  \State $\mathcal{S} \leftarrow \emptyset$
  \For{mỗi $\mathbf{x}_i \in \mathcal{D}_{\min}$}
    \State Tính $k$ láng giềng gần nhất của $\mathbf{x}_i$ trong $\mathcal{D}_{\min}$
  \EndFor
  \For{$t = 1$ \textbf{to} $N$}
    \State Chọn ngẫu nhiên $\mathbf{x}_i \in \mathcal{D}_{\min}$
    \State Chọn ngẫu nhiên $\mathbf{x}_{nn}$ từ $k$ láng giềng của $\mathbf{x}_i$
    \State $\lambda \sim \mathrm{Uniform}(0, 1)$
    \State $\mathbf{x}_{\text{new}} \leftarrow \mathbf{x}_i + \lambda(\mathbf{x}_{nn} - \mathbf{x}_i)$
    \State $\mathcal{S} \leftarrow \mathcal{S} \cup \{\mathbf{x}_{\text{new}}\}$
  \EndFor
  \State \Return $\mathcal{S}$
\end{algorithmic}
\end{algorithm}

\textbf{Ưu điểm:} Tạo ra dữ liệu tổng hợp có tính đa dạng cao hơn ROS, giảm nguy cơ overfitting, và giúp mô hình học được ranh giới quyết định tổng quát hơn.

\textbf{Nhược điểm:} SMOTE tạo mẫu tổng hợp một cách đồng đều trong toàn bộ không gian đặc trưng của lớp thiểu số mà không quan tâm đến mức độ khó phân loại của từng vùng. Điều này có thể dẫn đến việc tạo mẫu trong vùng chồng chéo giữa các lớp (\textit{overlapping region}), làm tăng nhiễu. Ngoài ra, SMOTE không hiệu quả khi dữ liệu có chiều cao (\textit{high-dimensional data}).

\textbf{Trường hợp nên sử dụng:} Phù hợp với tập dữ liệu có kích thước vừa, dữ liệu dạng số liên tục, và khi các lớp có vùng chồng chéo không quá lớn.

% ------------------------------------------------------------
\subsubsection{ADASYN (Adaptive Synthetic Sampling)}

\textbf{Khái niệm.}
ADASYN, được đề xuất bởi He et al. (2008), là một cải tiến của SMOTE theo hướng thích nghi (\textit{adaptive}). Thay vì phân bổ số lượng mẫu tổng hợp đồng đều, ADASYN tập trung tạo nhiều mẫu hơn ở những vùng khó phân loại --- tức là những mẫu thiểu số có nhiều láng giềng thuộc lớp đa số.

\textbf{Nguyên lý hoạt động.}
Với mỗi mẫu $\mathbf{x}_i$ trong lớp thiểu số, ADASYN tính tỉ lệ láng giềng thuộc lớp đa số trong $k$ láng giềng gần nhất:

\begin{equation}
  r_i = \frac{\Delta_i}{k},
  \quad \Delta_i = \text{số láng giềng thuộc lớp đa số}
  \label{eq:adasyn_r}
\end{equation}

Sau khi chuẩn hóa $\hat{r}_i = r_i \,/\, \sum_{j} r_j$, số lượng mẫu tổng hợp cần tạo cho mỗi $\mathbf{x}_i$ là:

\begin{equation}
  g_i = \hat{r}_i \cdot G
  \label{eq:adasyn_g}
\end{equation}

\noindent trong đó $G$ là tổng số mẫu cần tạo thêm. Mẫu tổng hợp được tạo ra theo cách tương tự SMOTE nhưng với mật độ phân bổ không đồng đều, tập trung vào vùng biên giới khó phân loại.

\begin{algorithm}[H]
\caption{ADASYN}
\begin{algorithmic}[1]
  \Require $\mathcal{D}_{\min}$, $\mathcal{D}_{\text{maj}}$, số láng giềng $k$, tỉ lệ cân bằng $d$
  \Ensure Tập mẫu tổng hợp $\mathcal{S}$
  \State $G \leftarrow (|\mathcal{D}_{\text{maj}}| - |\mathcal{D}_{\min}|) \times d$
  \For{mỗi $\mathbf{x}_i \in \mathcal{D}_{\min}$}
    \State Tìm $k$ láng giềng gần nhất; tính $r_i = \Delta_i / k$
  \EndFor
  \State Chuẩn hóa: $\hat{r}_i \leftarrow r_i / \sum_j r_j$
  \For{mỗi $\mathbf{x}_i \in \mathcal{D}_{\min}$}
    \State $g_i \leftarrow \lfloor \hat{r}_i \cdot G \rfloor$
    \For{$t = 1$ \textbf{to} $g_i$}
      \State Chọn ngẫu nhiên $\mathbf{x}_{nn}$ từ $k$ láng giềng của $\mathbf{x}_i$
      \State $\lambda \sim \mathrm{Uniform}(0,1)$
      \State $\mathbf{x}_{\text{new}} \leftarrow \mathbf{x}_i + \lambda(\mathbf{x}_{nn} - \mathbf{x}_i)$
      \State $\mathcal{S} \leftarrow \mathcal{S} \cup \{\mathbf{x}_{\text{new}}\}$
    \EndFor
  \EndFor
  \State \Return $\mathcal{S}$
\end{algorithmic}
\end{algorithm}

\textbf{Ưu điểm:} Thích nghi với cấu trúc phức tạp của dữ liệu, tập trung cải thiện hiệu năng ở vùng khó phân loại, và thường cho kết quả tốt hơn SMOTE khi ranh giới lớp phức tạp.

\textbf{Nhược điểm:} Phức tạp hơn và tốn kém hơn về mặt tính toán so với SMOTE. Việc tập trung vào vùng biên giới cũng có thể làm tăng nhiễu nếu dữ liệu gốc đã có nhiều nhiễu.

\textbf{Trường hợp nên sử dụng:} Hiệu quả khi ranh giới quyết định giữa các lớp phức tạp và mô hình gặp khó khăn đặc biệt ở vùng chuyển tiếp giữa hai lớp.

% ------------------------------------------------------------
\subsection{Undersampling Methods}
% ------------------------------------------------------------

\subsubsection{Random Undersampling (RUS)}

\textbf{Khái niệm và nguyên lý.}
Random Undersampling là kỹ thuật đơn giản nhất trong nhóm undersampling, loại bỏ ngẫu nhiên các mẫu từ lớp đa số cho đến khi đạt được tỉ lệ lớp mong muốn. Quá trình này giảm kích thước tổng thể của tập huấn luyện, giúp giảm độ lệch của mô hình về phía lớp đa số.

\begin{algorithm}[H]
\caption{Random Undersampling (RUS)}
\begin{algorithmic}[1]
  \Require $\mathcal{D}_{\text{maj}}$, $\mathcal{D}_{\min}$, tỉ lệ mục tiêu $r$
  \Ensure $\mathcal{D}'_{\text{maj}}$
  \State $N_{\text{keep}} \leftarrow r \cdot |\mathcal{D}_{\min}|$
  \State $\mathcal{D}'_{\text{maj}} \leftarrow$ lấy mẫu ngẫu nhiên $N_{\text{keep}}$ phần tử từ $\mathcal{D}_{\text{maj}}$
  \State \Return $\mathcal{D}'_{\text{maj}} \cup \mathcal{D}_{\min}$
\end{algorithmic}
\end{algorithm}

\textbf{Ưu điểm:} Giảm thời gian huấn luyện đáng kể do thu nhỏ kích thước tập dữ liệu, và hiệu quả đặc biệt trên các tập dữ liệu rất lớn.

\textbf{Nhược điểm:} Rủi ro lớn nhất là mất thông tin (\textit{information loss}) --- các mẫu bị loại bỏ ngẫu nhiên có thể chứa đựng các đặc trưng quan trọng không được đại diện bởi các mẫu còn lại, dẫn đến mô hình kém tổng quát hóa.

\textbf{Trường hợp nên sử dụng:} Phù hợp khi tập dữ liệu rất lớn và việc mất một phần dữ liệu lớp đa số không làm giảm đáng kể thông tin học được, hoặc khi tài nguyên tính toán là yếu tố hạn chế.

% ------------------------------------------------------------
\subsubsection{Tomek Links}

\textbf{Khái niệm.}
Tomek Links là một kỹ thuật undersampling có chọn lọc, được đề xuất bởi Tomek (1976). Một cặp mẫu $(\mathbf{x}_i, \mathbf{x}_j)$ thuộc hai lớp khác nhau được gọi là một \textit{Tomek Link} nếu không tồn tại mẫu $\mathbf{x}_k$ nào thỏa mãn:

\begin{equation}
  d(\mathbf{x}_i, \mathbf{x}_k) < d(\mathbf{x}_i, \mathbf{x}_j)
  \quad \text{hoặc} \quad
  d(\mathbf{x}_j, \mathbf{x}_k) < d(\mathbf{x}_i, \mathbf{x}_j)
  \label{eq:tomek}
\end{equation}

\noindent tức là hai mẫu là láng giềng gần nhất của nhau nhưng thuộc lớp khác nhau.

\textbf{Nguyên lý hoạt động.}
Kỹ thuật này xác định tất cả các Tomek Links trong tập dữ liệu, sau đó loại bỏ mẫu thuộc lớp đa số trong mỗi cặp. Các mẫu này thường là các mẫu nằm gần ranh giới quyết định hoặc là nhiễu, do đó việc loại bỏ chúng giúp làm sạch ranh giới phân lớp.

\textbf{Ưu điểm:} Có chọn lọc hơn RUS --- chỉ loại bỏ những mẫu nằm gần biên giới và có khả năng gây nhầm lẫn cho mô hình, từ đó cải thiện chất lượng ranh giới quyết định.

\textbf{Nhược điểm:} Thường không đủ để cân bằng dữ liệu khi tỉ lệ mất cân bằng cao, và thường được sử dụng như bước tiền xử lý kết hợp với các kỹ thuật khác hơn là đứng độc lập.

\textbf{Trường hợp nên sử dụng:} Phù hợp như bước làm sạch dữ liệu sau oversampling, hoặc khi bài toán có nhiễu cao tại vùng biên giới giữa các lớp.

% ------------------------------------------------------------
\subsubsection{SMOTE-ENN}

\textbf{Khái niệm.}
SMOTE-ENN là một phương pháp lai (\textit{hybrid}), kết hợp SMOTE với quy tắc làm sạch \textit{Edited Nearest Neighbor} (ENN). Phương pháp này trước tiên áp dụng SMOTE để oversampling lớp thiểu số, sau đó áp dụng ENN để loại bỏ các mẫu từ cả hai lớp mà bị phân loại sai bởi $k$ láng giềng gần nhất của chúng.

\textbf{Nguyên lý hoạt động.}
ENN loại bỏ một mẫu $\mathbf{x}_i$ nếu đa số $k$ láng giềng gần nhất của nó thuộc lớp khác với nhãn của $\mathbf{x}_i$:

\begin{equation}
  \text{Loại bỏ } \mathbf{x}_i \text{ nếu }
  \sum_{j \in \mathrm{kNN}(\mathbf{x}_i)} \mathbf{1}[y_j \neq y_i]
  > \frac{k}{2}
  \label{eq:enn}
\end{equation}

\noindent Quy tắc này có tác dụng loại bỏ các mẫu nhiễu và mẫu nằm trong vùng chồng chéo, làm sạch ranh giới quyết định sau bước oversampling.

\begin{algorithm}[H]
\caption{SMOTE-ENN}
\begin{algorithmic}[1]
  \Require $\mathcal{D}$, số láng giềng $k$, tỉ lệ mục tiêu $r$
  \Ensure Tập dữ liệu cân bằng và làm sạch $\mathcal{D}''$
  \State $\mathcal{D}' \leftarrow \text{SMOTE}(\mathcal{D},\, k,\, r)$
  \For{mỗi $\mathbf{x}_i \in \mathcal{D}'$}
    \State Tìm $k$ láng giềng gần nhất của $\mathbf{x}_i$ trong $\mathcal{D}'$
    \If{đa số láng giềng có nhãn $\neq y_i$}
      \State $\mathcal{D}' \leftarrow \mathcal{D}' \setminus \{\mathbf{x}_i\}$
    \EndIf
  \EndFor
  \State \Return $\mathcal{D}''= \mathcal{D}'$
\end{algorithmic}
\end{algorithm}

\textbf{Ưu điểm:} Kết hợp được khả năng tăng cường lớp thiểu số của SMOTE với khả năng làm sạch dữ liệu của ENN, tạo ra tập dữ liệu cân bằng hơn và sạch hơn so với chỉ dùng SMOTE đơn thuần.

\textbf{Nhược điểm:} Phức tạp hơn về mặt tính toán, và việc loại bỏ nhiều mẫu qua ENN có thể làm mất một phần thông tin quan trọng trong một số trường hợp.

\textbf{Trường hợp nên sử dụng:} Hiệu quả khi dữ liệu có nhiều nhiễu và vùng chồng chéo lớn giữa các lớp, đặc biệt khi SMOTE đơn thuần cho kết quả không ổn định.

% ============================================================
\section{So Sánh Oversampling và Undersampling}
% ============================================================

Oversampling và undersampling đại diện cho hai triết lý đối lập trong xử lý mất cân bằng. Oversampling bảo toàn toàn bộ thông tin trong tập dữ liệu gốc nhưng tăng kích thước tập huấn luyện, dẫn đến thời gian huấn luyện dài hơn và nguy cơ overfitting nếu không sử dụng kỹ thuật tạo dữ liệu tổng hợp. Undersampling giảm thời gian huấn luyện nhưng chấp nhận rủi ro mất thông tin từ các mẫu bị loại bỏ. Về mặt thực tiễn, oversampling thường được ưa chuộng hơn khi tập dữ liệu có kích thước vừa và nhỏ, trong khi undersampling phù hợp hơn cho các tập dữ liệu rất lớn nơi thông tin dư thừa ở lớp đa số là đáng kể.

\begin{table}[h]
\centering
\caption{So sánh Oversampling và Undersampling}
\label{tab:over_under}
\renewcommand{\arraystretch}{1.35}
\begin{tabularx}{\textwidth}{lXX}
\toprule
\textbf{Tiêu chí} & \textbf{Oversampling} & \textbf{Undersampling} \\
\midrule
Mất thông tin     & Không                 & Có thể xảy ra \\
Kích thước dữ liệu & Tăng                 & Giảm \\
Thời gian huấn luyện & Dài hơn            & Ngắn hơn \\
Nguy cơ overfitting  & Cao hơn (ROS)      & Thấp hơn \\
Phù hợp với         & Dữ liệu nhỏ--vừa   & Dữ liệu rất lớn \\
\bottomrule
\end{tabularx}
\end{table}

% ============================================================
\section{So Sánh SMOTE và ADASYN}
% ============================================================

Cả SMOTE và ADASYN đều tạo ra dữ liệu tổng hợp thông qua nội suy kNN, nhưng khác nhau về chiến lược phân bổ. SMOTE phân bổ mẫu tổng hợp đồng đều trên toàn bộ không gian lớp thiểu số, giúp tăng cường đa dạng tổng quát nhưng có thể tạo mẫu ở những vùng không cần thiết. ADASYN tập trung vào các vùng khó phân loại bằng cách ưu tiên tạo mẫu ở những điểm thiểu số gần ranh giới với lớp đa số, cho phép mô hình tập trung học ở vùng quyết định phức tạp hơn. Trong thực nghiệm, ADASYN thường cho kết quả tốt hơn SMOTE khi ranh giới quyết định phức tạp, nhưng có thể kém ổn định hơn khi dữ liệu nhiễu cao do xu hướng tạo mẫu ở vùng biên giới đã nhiễu.

% ============================================================
\section{Vai Trò Trong Cải Thiện Hiệu Suất Mô Hình}
% ============================================================

Các kỹ thuật xử lý mất cân bằng đóng vai trò then chốt trong việc đưa mô hình về trạng thái học tập không thiên lệch. Bằng cách điều chỉnh phân phối lớp trong tập huấn luyện, các kỹ thuật này giúp mô hình phân bổ ``sự chú ý'' học tập cân bằng hơn giữa các lớp, dẫn đến ranh giới quyết định phản ánh chính xác hơn cấu trúc thực của dữ liệu. Điều này đặc biệt quan trọng đối với các mô hình như SVM và hồi quy logistic, vốn nhạy cảm với sự mất cân bằng phân phối lớp. Nhiều nghiên cứu thực nghiệm đã chỉ ra rằng kết hợp kỹ thuật xử lý mất cân bằng với các mô hình phân loại truyền thống thường cho kết quả cải thiện đáng kể về F1-Score và PR-AUC, trong khi accuracy tổng thể có thể giảm nhẹ do mô hình không còn thiên vị lớp đa số.

% ============================================================
\section{Ảnh Hưởng Đến Overfitting}
% ============================================================

Mối quan hệ giữa xử lý mất cân bằng và overfitting cần được xem xét cẩn thận. Random Oversampling, do nhân bản mẫu hiện có, trực tiếp làm tăng nguy cơ overfitting bằng cách cho phép mô hình ``ghi nhớ'' các mẫu được lặp lại. Ngược lại, SMOTE và ADASYN tạo ra dữ liệu tổng hợp mới trong không gian đặc trưng, giúp tăng tính đa dạng và giảm nguy cơ overfitting so với ROS. Tuy nhiên, nếu tỉ lệ oversampling quá cao hoặc dữ liệu gốc quá ít, ngay cả SMOTE cũng có thể dẫn đến overfitting. Để kiểm soát vấn đề này, điều quan trọng là chỉ áp dụng kỹ thuật xử lý mất cân bằng trên tập huấn luyện, tuyệt đối không áp dụng trên tập kiểm tra hay kiểm định --- một sai lầm phổ biến gọi là \textit{data leakage}.

% ============================================================
\section{Ứng Dụng Thực Tế}
% ============================================================

Trong lĩnh vực \textbf{y tế}, mất cân bằng dữ liệu là đặc trưng phổ biến: bệnh nhân mắc bệnh hiếm gặp luôn ít hơn bệnh nhân bình thường. Các kỹ thuật như SMOTE và ADASYN được áp dụng rộng rãi trong chẩn đoán ung thư, phát hiện bệnh tim và phân tích dữ liệu gen, nơi việc bỏ sót dương tính thật (âm tính giả) có thể có hậu quả nghiêm trọng.

Trong lĩnh vực \textbf{tài chính và phát hiện gian lận}, các giao dịch gian lận chiếm tỉ lệ rất nhỏ trong tổng số giao dịch. Các hệ thống phát hiện gian lận thẻ tín dụng, gian lận bảo hiểm hay rửa tiền đều dựa nặng vào các kỹ thuật xử lý mất cân bằng, kết hợp với các chỉ số đánh giá như PR-AUC và F1-Score để đảm bảo tỉ lệ phát hiện gian lận cao trong khi kiểm soát báo động giả.

Trong lĩnh vực \textbf{phân loại văn bản và sự kiện hiếm}, các bài toán như phát hiện thư rác, phân loại sự cố mạng hay nhận diện nội dung độc hại thường đối mặt với tỉ lệ mất cân bằng đáng kể, và các kỹ thuật hybrid như SMOTE-ENN thường được sử dụng để cải thiện hiệu suất phát hiện.

% ============================================================
\section{Điều Kiện và Trường Hợp Áp Dụng}
% ============================================================

Việc lựa chọn kỹ thuật xử lý mất cân bằng phụ thuộc vào nhiều yếu tố. Khi tập dữ liệu nhỏ và tỉ lệ mất cân bằng vừa phải, SMOTE là lựa chọn an toàn và hiệu quả. Khi ranh giới quyết định phức tạp và có nhiều vùng chồng chéo, ADASYN hoặc SMOTE-ENN nên được ưu tiên. Khi tập dữ liệu rất lớn và tài nguyên tính toán hạn chế, RUS hoặc Tomek Links có thể là lựa chọn thực tế hơn. Trong mọi trường hợp, không có một phương pháp duy nhất tốt nhất cho tất cả tình huống --- thực nghiệm có kiểm định chéo là bước không thể bỏ qua để xác định chiến lược tối ưu cho từng bài toán cụ thể.

\begin{table}[h]
\centering
\caption{Hướng dẫn lựa chọn kỹ thuật xử lý mất cân bằng}
\label{tab:guide}
\renewcommand{\arraystretch}{1.35}
\begin{tabularx}{\textwidth}{lX}
\toprule
\textbf{Tình huống} & \textbf{Kỹ thuật khuyến nghị} \\
\midrule
Dữ liệu nhỏ, IR vừa phải                  & SMOTE \\
Ranh giới phức tạp, chồng chéo lớn        & ADASYN, SMOTE-ENN \\
Dữ liệu rất lớn, tài nguyên hạn chế       & RUS, Tomek Links \\
Dữ liệu có nhiều nhiễu vùng biên giới     & Tomek Links, SMOTE-ENN \\
Cần baseline nhanh                         & ROS \\
IR rất cao ($> 100:1$)                     & ADASYN + Ensemble \\
\bottomrule
\end{tabularx}
\end{table}

% ============================================================
\section{Kết Luận}
% ============================================================

Xử lý mất cân bằng dữ liệu là một bước tiền xử lý không thể thiếu trong quy trình xây dựng mô hình học máy cho các bài toán thực tế. Sự phong phú của các kỹ thuật hiện có --- từ ROS đơn giản đến ADASYN thích nghi và SMOTE-ENN lai ghép --- phản ánh sự đa dạng và phức tạp của vấn đề trong các miền ứng dụng khác nhau. Điều quan trọng cần nhấn mạnh là xử lý mất cân bằng không phải là giải pháp vạn năng, mà cần được tích hợp một cách có chủ đích vào toàn bộ quy trình học máy, bao gồm lựa chọn đặc trưng, thiết kế mô hình và đánh giá hiệu suất bằng các chỉ số phù hợp. Trong bối cảnh của luận văn này, các kỹ thuật xử lý mất cân bằng đóng vai trò nền tảng phương pháp luận, đảm bảo rằng các mô hình phân loại được xây dựng và đánh giá một cách công bằng và có ý nghĩa thực tiễn cao.
